\input{../tex-inputs/latex-korrekturansicht-vorspann}

\section[Arthur und Olga Schnitzler an Gustav Schwarzkopf, 4. 5. 1904]{L04358 Arthur und Olga Schnitzler an Gustav Schwarzkopf, 4. 5. 1904}
\nopagebreak\mylabel{L04358v}
\rehead{ }\normalsize\beginnumbering\briefempfaengerindex{Schwarzkopf, Gustav@\textsc{Schwarzkopf, Gustav}!zzzSchnitzler, Olga@\emph{von Olga Schnitzler}!1904-05-041@{4. 5. 1904}|(be}\briefempfaengerindex{Schwarzkopf, Gustav@\textsc{Schwarzkopf, Gustav}!zzzSchnitzler, Arthur@\emph{von Arthur Schnitzler}!1904-05-041@{4. 5. 1904}|(be}
\toendnotes[C]{\smallbreak\pagebreak[2]}
\correspDesc{Versand  durch Arthur Schnitzler, Olga Schnitzler am 4. 5. 1904 in Rom
\newline{}Erhalt  durch Gustav Schwarzkopf am 6. 5. 1904 in Wien}\toendnotes[C]{\smallbreak}
\Standort{DLA, A:Schnitzler, HS.1985.1.1897.}
\physDesc{Bildpostkarte, 186 Zeichen
\newline{}Handschrift Arthur Schnitzler: Bleistift, deutsche Kurrent
\newline{}Handschrift Olga Schnitzler: Bleistift, lateinische Kurrent
\newline{}Versand: 1) Stempel: »\nobreak{}\oindex{Rom@\textbf{Rom}, \emph{Hauptstadt}|pwk}Roma Ferrovia\nobreak{}«.   2) Stempel: »\nobreak{}\oindex{I., Innere Stadt@\textbf{I., Innere Stadt}, \emph{Verwaltungsgebiet}|pwk}1/1 Wien 1, 6. 5. 0\textcolor{gray}{4}, 8–9½V, \textcolor{gray}{Bestellt}\nobreak{}«. }\pstart{}{\pb}\textcolor{pink}{Austria}\oindex{Österreich@\textbf{Österreich}|pw}{}\ledrightnote{\textcolor{pink}{Österreich}}\pend{}\pstart{}\textcolor{gray}{\textbf{A}}\textsc{l Signor Gustav
                     Schwarzkopf}\pend{}\pstart{}\textcolor{pink}{\textcolor{gray}{\textbf{Vienna}}}\oindex{Wien@\textbf{Wien}, \emph{Verwaltungsgebiet}|pw}{}\ledrightnote{\textcolor{pink}{Wien}}\pend{}\pstart{}\textsc{\textcolor{pink}{I. Tiefer
                        Graben 23}\oindex{Wien@\textbf{Wien}!I., Innere Stadt@\textbf{I., Innere Stadt}!Tiefer Graben 23@\textbf{Tiefer Graben 23}, \emph{Wohngebäude}|pw}{}\ledrightnote{\textcolor{pink}{Tiefer Graben 23}}.}\pend{}{\bigskip}
\pstart
           \noindent{}\centering{}{\pb}\textcolor{gray}{\textbf{\textcolor{pink}{Roma}\oindex{Rom@\textbf{Rom}, \emph{Hauptstadt}|pw}{}\ledrightnote{\textcolor{pink}{Rom}} – \textcolor{pink}{Foro
                     Romano}\oindex{Forum Romanum@\textbf{Forum Romanum}, \emph{Ausgrabung}|pw}{}\ledrightnote{\textcolor{pink}{Forum Romanum}} – \textcolor{pink}{Tempio di Castore e
                     Polluce}\oindex{Tempel von Castor und Pollux@\textbf{Tempel von Castor und Pollux}|pw}{}\ledrightnote{\textcolor{pink}{Tempel von Castor und Pollux}} e \textcolor{pink}{Basilica Giulia}\oindex{Basilica Julia@\textbf{Basilica Julia}|pw}{}\ledrightnote{\textcolor{pink}{Basilica Julia}}}}\pend
           \vspace{1em}
\pstart
           {\pb}Herzlichen Gruß{\\[\baselineskip]}Ihr \spacefill\mbox{A. S.}\pend
           \leftskip=0em{}
\pstart
           4. 5. 904\pend
           \selectlanguage{ngerman}\vspace{1em}
\pstart
           \noindent{}{[}hs. Schnitzler:{]} Woher weiß man, ob die alten
                  \textcolor{gray}{Älteren}{ }\textcolor{gray}{×}\-\textcolor{gray}{×}\-\textcolor{gray}{×}\-\textcolor{gray}{×}\-\textcolor{gray}{×}{ }\textcolor{gray}{ſin}d? ſagte S. u. machte ein ironiſches Geſicht.\pend
           
\pstart
           Herzl Grüße{\\[\baselineskip]}\spacefill\mbox{Olga}\pend
           \leftskip=0em{}\selectlanguage{ngerman}\endnumbering\briefempfaengerindex{Schwarzkopf, Gustav@\textsc{Schwarzkopf, Gustav}!zzzSchnitzler, Olga@\emph{von Olga Schnitzler}!1904-05-041@{4. 5. 1904}|)be}\briefempfaengerindex{Schwarzkopf, Gustav@\textsc{Schwarzkopf, Gustav}!zzzSchnitzler, Arthur@\emph{von Arthur Schnitzler}!1904-05-041@{4. 5. 1904}|)be}\mylabel{L04358h}
\begin{anhang}
\end{anhang}\newcommand{\dateiname}{L04358}\newcommand{\titel}{Arthur und Olga Schnitzler an Gustav Schwarzkopf, 4. 5. 1904}\newcommand{\editorInnen}{Herausgegeben von Jahnke, SelmaMüller, Martin Anton}\input{../tex-inputs/latex-korrekturansicht-abspann}
